\documentclass{article}

% ICLR 2027 style
\usepackage{iclr2027_conference}

\usepackage{times}
\usepackage{amsmath,amssymb,amsthm}
\usepackage{booktabs}
\usepackage{graphicx}
\usepackage{hyperref}
\usepackage{natbib}
\usepackage{multirow}
\usepackage{xcolor}

\newtheorem{proposition}{Proposition}
\newtheorem{remark}{Remark}

\title{Beyond High Prices: Auditing Information-Mediated Enforcement in Recurrent Pricing Agents}

\author{Anonymous}

\begin{document}

\maketitle

\begin{abstract}
Observing high prices or correlated price movements does not prove that pricing algorithms have learned collusive punishment. We study how to separate genuine dynamic enforcement from static pricing incentives in a continuous-price duopoly where firms observe only noisy, delayed public reports and learn via recurrent proximal policy optimization (PPO). First, we establish two analytical benchmarks: a two-state public trigger model characterizing when future continuation losses deter unilateral price cuts, and a fixed-threshold Gaussian monitoring model showing how observation noise weakens deterrence by raising both false alarms and missed detections. Second, we audit frozen GRU-PPO policies using paired counterfactual price-cut interventions and an action-independent signal ablation that serves as a proved negative control. Across ten training seeds and one million market interactions per run, normalized price elevation falls from $0.43$ in early training to $0.015$ under informative monitoring and $0.019$ after signal ablation, with a paired difference of $0.0037$ (95\% bootstrap CI: $[-0.0241,\,0.0288]$). Crucially, the payoff losses from forced price cuts occur almost entirely in the intervention period itself rather than through subsequent continuation losses, and several small deviations remain profitable. Rival price responses are negligible under informative monitoring and identically zero under ablation. Together, these findings distinguish apparent coordination from information-mediated enforcement and show no evidence of learned trigger punishment in this setting.
\end{abstract}

%% ============================================================
\section{Introduction}
%% ============================================================

High prices alone do not reveal why they persist. When pricing algorithms charge above the competitive benchmark, they may do so because exploration is incomplete, because function approximation introduces upward bias, or because unilateral price cuts trigger costly future retaliation. Only the third mechanism constitutes dynamic collusive enforcement. A credible audit of algorithmic pricing must therefore separate immediate static payoff effects from future continuation penalties rather than relying on converged price levels alone.

Our economic premise is straightforward: holding demand and cost primitives fixed, the public information channel governs whether coordinated prices can be enforced through continuation threats. In repeated games with imperfect public monitoring, the deterrent power of a trigger strategy depends on the gap between detecting a true deviation and falsely punishing cooperation, scaled by the continuation value at risk. We formalize this relationship in a fixed public trigger benchmark. Importantly, this benchmark is a conditional characterization of dynamic incentives; it does not assume that gradient-based neural learners will converge to such a strategy.

Our empirical question is correspondingly focused: do recurrent PPO pricing agents learn policies whose price levels and responses to unilateral price cuts depend on an action-sensitive public signal? To answer this, we compare an informative public reporting channel against an action-independent replacement channel, retraining both firms from scratch under identical budgets. We then freeze the trained policies and inject paired one-period price cuts to decompose total payoff changes into immediate impact effects and discounted continuation effects.

In the tested price-cut probes, we find no evidence of information-mediated enforcement by the recurrent PPO agents. Average prices converge near the one-shot Nash benchmark under both informative and action-independent channels, and removing action information produces no statistically resolved drop in final prices or gain in the archived welfare index. Furthermore, counterfactual price cuts lose money immediately at impact rather than through future retaliation, while action-independent monitoring yields an identically zero impulse response by architectural invariance. Our contribution is an integrated framework combining (i) a corrected analytical benchmark for noisy public triggers, (ii) a counterfactual policy audit paired with a proved negative control, and (iii) a complete ten-seed recurrent PPO evaluation.

\subsection{Research Gap: Imperfect Public Monitoring and Recurrent Policy Auditing}

Classical repeated-game theory shows how firms can sustain supra-competitive prices under imperfect public monitoring when rival actions are hidden and only noisy public signals are observed \citep{green1984noncooperative, abreu1990toward, fudenberg1994folk}. In that literature, deviations are deterred when adverse public signals shift play into lower-value continuation regimes, balancing deviation detection against costly false alarms. In algorithmic pricing, \citet{calvano2020artificial} show that tabular Q-learning sustains supra-competitive prices under perfect price observation, and \citet{calvano2021algorithmic} extend tabular Q-learning to imperfect monitoring adapted from \citet{green1984noncooperative}, studying quantity-setting Cournot learners with stochastic demand, public market-price signals, and finite-duration ``price wars''. Meanwhile, recent deep reinforcement learning studies such as \citet{deng2025exploring} compare tabular Q-learning, DQN, and PPO under perfect monitoring, focusing primarily on converged price levels.

Two gaps separate these prior studies from our work:
\begin{enumerate}
\item \textbf{Difference from prior imperfect public monitoring studies:} Classical theory \citep{green1984noncooperative, abreu1990toward, fudenberg1994folk} studies public-strategy incentives in infinite-horizon games, while \citet{calvano2021algorithmic} also use forced output deviations to examine punishment paths. However, observing high prices or post-shock price dips does not by itself prove that future continuation losses outweigh immediate deviation gains. We extend these diagnostics to continuous-price neural policies under delayed, two-period-averaged Gaussian reports by decomposing counterfactual price-cut payoffs into \emph{immediate impact gain} versus \emph{discounted continuation value change}, paired with an exact action-independent channel replacement (Proposition~\ref{prop:invariance}) as a proved zero-response negative control. We do not claim to refute \citet{calvano2021algorithmic} or generalize across all learners; rather, we test whether observed policy responses reflect dynamic continuation deterrence or immediate static losses.
\item \textbf{Why a windowed GRU with PPO is used:} In our market, firms observe neither rival prices nor contemporaneous one-period signals. Instead, each public report averages two consecutive latent price states, adds Gaussian noise, and arrives after a one-step queue delay. Because of this aggregation, delay, and noise---the signal-extraction problem of imperfect monitoring \citep{green1984noncooperative, calvano2021algorithmic}---a single scalar report cannot distinguish a transient noise shock from a unilateral price cut. Filtering noise requires conditioning on a sequence of past reports, while continuous prices in $[1,\,2.5]$ and continuous reports make tabular discretization over multi-period histories intractable. We therefore pair continuous-action PPO \citep{schulman2017proximal} with a 128-unit GRU encoding a fixed window of 32 public reports (recomputed from a fresh hidden state at each step) as a compact finite-memory sequence encoder. We use GRU+PPO strictly as a controlled testbed for history-dependent continuous pricing policies, not as a claim that it represents commercial pricing software or outperforms tabular, MLP, or no-regret learners (see Appendix~A.5 and~C).
\end{enumerate}

%% ============================================================
\section{Market and Observation Model}
%% ============================================================

Two symmetric firms $i\in\{1,2\}$ simultaneously choose continuous prices $p_{i,t}\in[1,\,2.5]$. Marginal cost is $c=1$, product quality is $v=2$, outside option utility is $0$, market size is $1$, and the logit scale parameter is $\mu=0.25$. For firm $i$ against rival $j$, logit demand, stage profit, and logit total surplus (with outside-option normalization) are
\begin{align}
q_i(p_i,p_j) &= \frac{\exp\!\bigl((v-p_i)/\mu\bigr)}
{1+\exp\!\bigl((v-p_i)/\mu\bigr)+\exp\!\bigl((v-p_j)/\mu\bigr)},\\[3pt]
\pi_i(p_i,p_j) &= (p_i - c)\,q_i(p_i,p_j),\\[3pt]
\mathrm{TS}(p_i,p_j) &= \mu\log\!\left(1+\sum_{k\in\{i,j\}}e^{(v-p_k)/\mu}\right)+\pi_i(p_i,p_j)+\pi_j(p_j,p_i).
\end{align}
The archived welfare index is retained as reported; its implementation cannot be checked without the simulator source (Appendix~G). The $1{,}501$-point grid references used for reported elevation are $p^{\mathrm{N}}=1.473$ and $p^{\mathrm{J}}=1.925$. Continuous first-order conditions give $1.472927$ and $1.924981$, respectively (Appendix~B.1). We measure normalized price elevation as
\[
E = \frac{\bar{p}-p^{\mathrm{N}}}{p^{\mathrm{J}}-p^{\mathrm{N}}},
\]
which records where average price $\bar{p}$ lies relative to the Nash ($E=0$) and joint-profit ($E=1$) benchmarks without clipping to $[0,1]$.

\subsection{Gaussian Monitoring and Channel Replacement}

At the start of each episode, a monitoring parameter vector $\theta=(\theta_{\mathrm{lo}},\theta_{\mathrm{hi}},\sigma)\in\Theta$ is drawn uniformly from a retained MCMC calibration posterior and held fixed for the episode (the stated calibration model and missing replication inputs are described in Appendix~B). Current prices determine the latent monitoring mean via the minimum price:
\begin{equation}\label{eq:latent}
m_t = \theta_{\mathrm{lo}}
+ (\theta_{\mathrm{hi}}-\theta_{\mathrm{lo}})
\cdot \frac{\min(p_{1,t},p_{2,t})-p^{\mathrm{N}}}
{p^{\mathrm{J}}-p^{\mathrm{N}}},
\end{equation}
clamped to $[\theta_{\mathrm{lo}},\theta_{\mathrm{hi}}]$. Under the \textbf{informative channel}, the public report averages the two most recent latent means and adds a Gaussian innovation:
\begin{equation}\label{eq:report}
r_{t+1} = \frac{m_t + m_{t-1}}{2} + \varepsilon_t,\qquad
\varepsilon_t\sim\mathcal{N}(0,\sigma_\varepsilon^2).
\end{equation}
Here $\sigma_\varepsilon$ denotes the effective report-noise scale; the configuration also records an added scale of $0.1$, whose combination with posterior $\sigma$ cannot be checked from the archive alone (Appendix~G). Reports pass through a one-period queue delay, so $r_{t+1}$ enters the observation window only after period-$(t+1)$ actions are chosen. Delivered reports are normalized by centering at $1.6$, dividing by $0.4$, and clipping to $[-5,5]$. Under the \textbf{action-independent ablation channel}, reports are drawn directly as $r_{t+1}\sim\mathcal{N}(1.6,\,0.4^2)$ independently of prices, keeping delay, observation dimensions, and market primitives unchanged. Because this replacement alters both action dependence and marginal report moments, comparing retrained policies identifies the effect of the full channel replacement rather than a pure information effect with matched moments (Appendix~B.2).

%% ============================================================
\section{Learning Architecture and Counterfactual Audit}
%% ============================================================

Each firm maintains an independent actor--critic network trained via PPO \citep{schulman2017proximal}. At each step, the actor observes the window of the 32 most recent normalized public reports and normalized episode time $t/512$, but observes neither rival prices nor realized profits during execution. A 128-unit GRU encodes the 32-report window from a fresh zero hidden state, and its final state is concatenated with time and passed through a linear projection, $\tanh$, and layer normalization. Each run trains for $10^6$ joint market interactions across $512$-period episodes with discount factor $\delta=0.99$ and no terminal value bootstrap. During training, actions are sampled from a $\tanh$-squashed Gaussian distribution; during evaluation and auditing, policies are frozen and deployed deterministically using the latent mean action $\mu(s)$ (complete hyperparameters and loss definitions are in Appendix~C).

\subsection{Counterfactual Price-Cut Interventions}

To test whether frozen policies deter price cuts through future punishment, we run paired counterfactual rollouts. For each trained policy pair, we intervene on firm $0$ or firm $1$ at period $t_0\in\{64,256\}$ by forcing a one-period normalized action cut of $0.1$ or $0.3$ (corresponding to unclipped price cuts of $0.075$ and $0.225$). This yields $2\times 2\times 2=8$ probe types per seed, each evaluated over 128 common-random-number rollouts sharing the same posterior draw $\theta$ and noise innovations $\varepsilon_t$. After period $t_0$, both firms revert to their frozen policies through period $511$. We track the rival impulse response
\[
\Delta p^{\mathrm{rival}}_t
= p^{\mathrm{rival}}_t(\text{cut branch})
- p^{\mathrm{rival}}_t(\text{control branch})
\]
and decompose the deviating firm's discounted payoff change into the immediate impact gain at $t_0$ and the continuation gain $\Delta\pi^{\mathrm{cont}}$ over $t>t_0$.

\subsection{An Exact Negative Control}

\begin{proposition}[Observation invariance]\label{prop:invariance}
Suppose execution policies are frozen deterministic functions of public report histories and time. If reports are action independent, initial buffers and innovations are coupled across branches, and actions do not enter any other actor input, then a one-period action intervention at $t_0$ leaves all subsequent prescribed actions unchanged. Consequently, $\Delta p^{\mathrm{rival}}_t = 0$ for all $t > t_0$ and $\Delta\pi^{\mathrm{cont}} = 0$ pathwise.
\end{proposition}

\begin{proof}
Prior to $t_0$, both branches share identical inputs. At $t_0$, the non-intervened rival observes the same history and selects the same action. Because the report generator under ablation does not depend on prices and shares coupled innovations, the newly queued and delivered reports remain identical across branches at $t_0+1$. By induction, all future observation windows and prescribed actions coincide for all $t>t_0$. Since stage payoffs depend only on current prices, every post-intervention payoff difference is zero, leaving only the impact-period difference at $t_0$. \qed
\end{proof}

%% ============================================================
\section{Experimental Evidence Across Training Seeds}
%% ============================================================

We train ten independent seed pairs from scratch under both the informative and action-independent channels. Each run logs 100 diagnostic checkpoints (16 episodes each) across training and conducts a final evaluation over 64 episodes (seed-by-seed outcomes are reported in Appendix~D).

\begin{figure}[t]
    \includegraphics[width=1.0\linewidth]{media/image1.png}
\caption{Deterministic deployment across training. Solid lines average ten training seeds; shaded bands show pointwise 95\% confidence intervals from 2{,}000 whole-seed bootstrap resamples. Zero elevation corresponds to the one-shot Nash price $p^{\mathrm{N}}$.}
\label{fig:elevation}
\end{figure}

As shown in Figure~\ref{fig:elevation}, normalized price elevation is near $0.43$ at the first recorded checkpoint, after $12{,}288$ training interactions, in both conditions ($0.4318$ informative vs.\ $0.4327$ action independent), and then declines overall toward the Nash benchmark with intermediate fluctuations. By the final evaluation (Table~\ref{tab:final}), mean price elevation is $0.0151$ under informative monitoring and $0.0188$ under action-independent monitoring, with a paired difference of $0.0037$ (95\% bootstrap CI: $[-0.0241,\,0.0288]$).

\begin{table}[t]
\centering
\caption{Final evaluation across ten seeds (64 episodes per seed). Differences are action independent minus informative, with pointwise 95\% paired bootstrap confidence intervals (2{,}000 resamples).}
\label{tab:final}
\begin{tabular}{@{}lrrr@{}}
\toprule
Outcome & Informative & Action independent & Paired diff.\ [95\% CI]\\
\midrule
Mean price
  & 1.479804 & 1.481492 & 0.001688 $[-0.0109,\,0.0130]$\\
Price elevation
  & 0.015052 & 0.018787 & 0.003735 $[-0.0241,\,0.0288]$\\
Archived welfare index
  & 1.159588 & 1.159660 & 0.000072 $[-0.0013,\,0.0017]$\\
Feature rank (firm 0)
  & 1.274335 & 1.338647 & 0.064312 $[0.0113,\,0.1337]$\\
Feature rank (firm 1)
  & 1.269538 & 1.353578 & 0.084040 $[0.0282,\,0.1495]$\\
\bottomrule
\end{tabular}
\end{table}

Removing action information from the public channel therefore produces no statistically resolved reduction in converged prices or improvement in the archived welfare index. The price-elevation interval spans zero and effects of either sign; without a prespecified equivalence margin, it does not establish equivalence or rule out all economically important effects. Effective feature rank (the exponential entropy of normalized eigenvalues of centered encoder features; Appendix~D) is slightly higher under ablation. This descriptive difference does not identify its cause or a causal pricing mechanism.

%% ============================================================
\section{Deviation Incentives and Rival Responses}
%% ============================================================

A price cut that reduces total profit is not necessarily deterred by rival retaliation; cutting prices near the Nash benchmark can simply be unprofitable immediately. Table~\ref{tab:audit} and Figure~\ref{fig:deviation} therefore compare immediate impact gains against full remaining-episode gains across the two cut sizes, treating the ten training seeds as the unit of replication (all 160 individual probe estimates appear in Appendix~F; uncertainty conventions are in Appendix~E).

\begin{figure}[t]
    \includegraphics[width=1.0\linewidth]{media/image_2.png}
\caption{Left: mean rival price response $\Delta p^{\mathrm{rival}}_t$ across eight probe types and ten seeds (in $10^{-6}$ price units) with 95\% seed-bootstrap intervals. Right: immediate impact gain versus full remaining-episode gain for each seed, averaged by cut size (circles: $0.1$ cut; squares: $0.3$ cut). Points lie on the $45^\circ$ equality line because continuation effects are negligible.}
\label{fig:deviation}
\end{figure}

\begin{table}[t]
\centering
\caption{Counterfactual price-cut audit across ten seeds. ``Positive probes'' counts positive conditional mean full gains among the 40 seed--firm--time probes per cut size and condition.}
\label{tab:audit}
\begin{tabular}{@{}llrrr@{}}
\toprule
Channel & Action cut & Immediate gain & Full gain & Positive probes\\
\midrule
Informative       & 0.1 & $-0.005186$ & $-0.005192$ & 3 of 40\\
Informative       & 0.3 & $-0.051317$ & $-0.051322$ & 0 of 40\\
Action independent & 0.1 & $-0.004800$ & $-0.004800$ & 1 of 40\\
Action independent & 0.3 & $-0.050273$ & $-0.050273$ & 0 of 40\\
\bottomrule
\end{tabular}
\end{table}

Under informative monitoring, the average full gain is $-0.005192$ for the $0.1$ action cut and $-0.051322$ for the $0.3$ action cut, while the discounted continuation component $\Delta\pi^{\mathrm{cont}}$ averages only $-5\times 10^{-6}$ in both cases. Thus, about $99.9\%$ or more of the mean deviation loss is incurred immediately at the intervention step $t_0$ because baseline prices are already near $p^{\mathrm{N}}$. Under action-independent monitoring, immediate and full gains are identical to machine precision ($\Delta\pi^{\mathrm{cont}}=0$), confirming Proposition~\ref{prop:invariance}.

Inspecting individual probes also reveals that not all deviations are unprofitable: 3 of 40 small-cut probes under informative monitoring (across two seeds, max gain $0.003846$) and 1 of 40 under ablation (gain $0.000122$) yield positive conditional mean gains. The signed rival-response area under informative monitoring averages $-1.36\times 10^{-5}$ price-period units (95\% seed-bootstrap CI: $[-2.53\times 10^{-5},\,-3.42\times 10^{-6}]$). This small downward response is statistically resolved, but its magnitude and the negligible continuation losses do not establish dynamic enforcement. We also note that because Equation~\eqref{eq:latent} clamps latent means below $p^{\mathrm{N}}$, cuts from near-Nash prices can saturate the signal, limiting sensitivity to shock size (Appendix~B.2).

%% ============================================================
\section{A Fixed Public Trigger Benchmark}
%% ============================================================

To clarify what dynamic enforcement would require, consider an infinite-horizon repeated stage game with discount factor $\delta\in[0,1)$, a cooperative public state $C$, and an absorbing punishment state $P$ prescribing credible stage-game Nash play ($\pi^P$). In state $C$, firms earn $\pi^C > \pi^P$. A bad public signal occurs with false-alarm probability $\alpha\in[0,1)$ under cooperation and detection probability $\beta\in(\alpha,1]$ after a specified one-period unilateral deviation yielding current profit $\pi^D > \pi^C$.

\begin{proposition}[Incentive threshold]\label{prop:threshold}
The specified one-period deviation is unprofitable in state $C$ if and only if
\begin{equation}\label{eq:threshold}
\delta \;\geq\;
\frac{\pi^D - \pi^C}
{(\pi^D-\pi^C)(1-\alpha) + (\beta-\alpha)(\pi^C-\pi^P)}.
\end{equation}
Equivalently, $\delta\geq\delta^*$, where
\[
\delta^* = \frac{g}{g(1-\alpha) + \Delta\cdot\Delta V},
\]
with $g=\pi^D-\pi^C>0$, $\Delta=\beta-\alpha>0$, and $\Delta V=\pi^C-\pi^P>0$. A feasible discount factor $\delta^*<1$ exists if and only if $\Delta\cdot\Delta V > \alpha g$.
\end{proposition}

\begin{proof}
The Bellman equations for cooperation, one-period deviation, and absorbing punishment are
\begin{align}
V^C &= \pi^C + \delta\bigl[(1-\alpha)V^C + \alpha V^P\bigr],
\label{eq:vc}\\
V^D &= \pi^D + \delta\bigl[(1-\beta)V^C + \beta V^P\bigr], \qquad
V^P = \pi^P + \delta V^P.
\label{eq:vd}
\end{align}
Subtracting~\eqref{eq:vc} from~\eqref{eq:vd} yields $V^D - V^C = g - \delta\Delta(V^C-V^P)$. Subtracting $V^P=\pi^P+\delta V^P$ from~\eqref{eq:vc} gives $[1-\delta(1-\alpha)](V^C-V^P)=\pi^C-\pi^P$, so
\[
V^C - V^P = \frac{\Delta V}{1-\delta(1-\alpha)} > 0.
\]
Substituting $V^C-V^P$ into $V^D-V^C\leq 0$ gives $g[1-\delta(1-\alpha)] \leq \delta\Delta\cdot\Delta V$. Rearranging for $\delta$ yields~\eqref{eq:threshold}, and requiring $\delta^*<1$ gives $\Delta\cdot\Delta V > \alpha g$ (extended derivations and numerical examples appear in Appendix~A.1--A.2). \qed
\end{proof}

If $\beta\leq\alpha$ and $g>0$, then $V^D-V^C=g-\delta(\beta-\alpha)(V^C-V^P)>0$, so no admissible discount factor deters that deviation. Full incentive compatibility in state $C$ requires deterring every available unilateral action $a_i$, not only the one maximizing current profit. With $g(a_i)=\pi_i(a_i,a_{-i}^C)-\pi^C$ and $\Delta(a_i)=\beta(a_i)-\alpha$, this requires
\begin{equation}\label{eq:full}
g(a_i)\bigl[1-\delta(1-\alpha)\bigr] \;\leq\;
\delta\,\Delta(a_i)\,\Delta V
\quad\text{for every firm } i \text{ and action } a_i.
\end{equation}
Because small price cuts may have low detection $\Delta(a_i)$, a finite set of probes cannot certify~\eqref{eq:full} over a continuous action space. Differentiating $\delta^*$ with $D=g(1-\alpha)+\Delta\cdot\Delta V > 0$ yields the comparative statics:
\begin{equation}\label{eq:compstat}
\frac{\partial\delta^*}{\partial\alpha}=\frac{g(g+\Delta V)}{D^2}>0,\qquad
\frac{\partial\delta^*}{\partial\beta}=-\frac{g\,\Delta V}{D^2}<0,\qquad
\frac{\partial\delta^*}{\partial(\pi^C-\pi^P)}=-\frac{g\Delta}{D^2}<0.
\end{equation}
Higher false alarms $\alpha$ erode cooperative continuation value $V^C-V^P$ and strictly raise the patience required for enforcement, while higher detection $\beta$ and larger punishment gaps $\Delta V$ lower it.

%% ============================================================
\section{What the Gaussian Assumption Proves}
%% ============================================================

Suppose the public signal is drawn from $Y_t\sim\mathcal{N}(\mu_C,\sigma^2)$ under cooperation and $Y_t\sim\mathcal{N}(\mu_D,\sigma^2)$ after the specified deviation, and punishment is triggered when $Y_t<\tau$ for a fixed threshold $\mu_D<\tau<\mu_C$. Then
\[
\alpha = \Phi\!\left(\frac{\tau-\mu_C}{\sigma}\right),\qquad
\beta  = \Phi\!\left(\frac{\tau-\mu_D}{\sigma}\right).
\]

\begin{proposition}[Restricted noise comparative static]\label{prop:gauss}
Holding $\mu_C$, $\mu_D$, $\tau$, prescribed actions, and payoff gaps fixed with $\mu_D<\tau<\mu_C$, increasing noise $\sigma>0$ strictly decreases $\Delta=\beta-\alpha$ and strictly raises $\delta^*$ in~\eqref{eq:threshold}.
\end{proposition}

\begin{proof}
Let $z_C=(\tau-\mu_C)/\sigma < 0$ and $z_D=(\tau-\mu_D)/\sigma > 0$. With standard normal density $\phi$,
\begin{equation}\label{eq:ddelta}
\frac{\partial\Delta}{\partial\sigma}
= -\frac{1}{\sigma}\bigl[z_D\,\phi(z_D)-z_C\,\phi(z_C)\bigr] < 0,
\end{equation}
since $\alpha'(\sigma)=-z_C\phi(z_C)/\sigma > 0$ and $\beta'(\sigma)=-z_D\phi(z_D)/\sigma < 0$. Combining~\eqref{eq:ddelta} with~\eqref{eq:compstat} yields $\partial\delta^*/\partial\sigma = g[(g+\Delta V)\alpha'(\sigma)-\Delta V\,\beta'(\sigma)]/D^2 > 0$. As $\sigma\to 0^+$, $\alpha\to 0$ and $\beta\to 1$, so $\delta^*\to g/(g+\Delta V)$. As $\sigma\to\infty$, both $\alpha,\beta\to 1/2$, so $\delta^*\to 2$, which is infeasible for $\delta<1$ (Appendix~A.3). \qed
\end{proof}

\subsection{Conditional Connection to the PPO Environment}

Propositions~\ref{prop:threshold} and~\ref{prop:gauss} provide an exact incentive accounting benchmark, but three structural differences separate this analytical model from the PPO experiment: (i) the PPO simulator delivers continuous, two-period-averaged, one-step-delayed reports to a 32-step GRU rather than enforcing a contemporaneous binary threshold; (ii) the experiment compares an informative channel against an action-independent replacement rather than sweeping $\sigma$ across a fixed trigger policy; and (iii) PPO episodes terminate at $T=512$ with observable normalized time rather than running over an infinite horizon. Consequently, our empirical audit tests whether trained recurrent policies exhibit continuation-value deterrence in practice, without claiming that finite-horizon PPO must replicate infinite-horizon trigger equilibria.

%% ============================================================
\section{Interpretation, Limitations, and Conclusion}
%% ============================================================

Imperfect public monitoring theory shows that sustaining collusion requires linking noisy public signals to continuation value losses \citep{green1984noncooperative, abreu1990toward, fudenberg1994folk}. While tabular Q-learning can learn supra-competitive pricing under both perfect monitoring \citep{calvano2020artificial} and imperfect monitoring with finite price wars \citep{calvano2021algorithmic}, and deep RL studies report near-Nash pricing under perfect-monitoring PPO \citep{deng2025exploring}, our contribution lies in auditing \emph{how} deviations are penalized via counterfactual payoff decomposition and channel replacement.

Five methodological limitations bound our conclusions (expanded in Appendix~G--H):
\begin{itemize}
\item \textbf{Input distribution confound:} Replacing the informative channel with $r_{t+1}\sim\mathcal{N}(1.6,0.4^2)$ alters marginal report moments alongside action dependence; isolating a pure information effect requires moment-matched placebo channels.
\item \textbf{Signal saturation and probe scope:} Clamping in~\eqref{eq:latent} saturates latent means at or below $p^{\mathrm{N}}$, and our audit tests a finite set of one-step cuts rather than best-response or multi-step deviations.
\item \textbf{Missing positive control:} Proposition~\ref{prop:invariance} validates our zero-response negative control, but evaluating audit sensitivity (statistical power) requires testing against a known collusive trigger positive control.
\item \textbf{Finite horizon and model scope:} Known $512$-period termination with time inputs and ten training seeds in a single stationary market bound generalizability; an interval spanning zero does not prove exact treatment equivalence.
\item \textbf{Reproducibility accounting:} The archive records 20 completed PPO jobs ($58.17$M interactions; $44.06$ hours) and MCMC diagnostics ($\hat{R}=1.008945$, ESS $\ge 606.86$ \citep{vehtari2021rank}). Checkpoint hashes verify file integrity; missing source and calibration inputs prevent independent training replication (Appendix~G).
\end{itemize}

\paragraph{Conclusion.}
High prices and post-shock price movements must not be conflated with dynamic collusive enforcement. By combining a corrected noisy-trigger benchmark with a counterfactual price-cut audit and an exact negative control, we show that in this recurrent PPO setting, early high prices decay to near-Nash levels and deviation losses are overwhelmingly immediate rather than continuation-driven. The auditor observes simulated prices, while the agents receive only public reports; information-mediated enforcement remains a structural hypothesis to test directly.

%% ============================================================
%% References
%% ============================================================
\bibliographystyle{iclr2027_conference}
\begin{thebibliography}{9}

\bibitem[Abreu et~al.(1990)]{abreu1990toward}
D.~Abreu, D.~Pearce, and E.~Stacchetti.
\newblock Toward a theory of discounted repeated games with imperfect
  monitoring.
\newblock \emph{Econometrica}, 58(5):1041--1063, 1990.

\bibitem[Calvano et~al.(2020)]{calvano2020artificial}
E.~Calvano, G.~Calzolari, V.~Denicol\`o, and S.~Pastorello.
\newblock Artificial intelligence, algorithmic pricing, and collusion.
\newblock \emph{American Economic Review}, 110(10):3267--3297, 2020.

\bibitem[Calvano et~al.(2021)]{calvano2021algorithmic}
E.~Calvano, G.~Calzolari, V.~Denicol\`o, and S.~Pastorello.
\newblock Algorithmic collusion with imperfect monitoring.
\newblock \emph{International Journal of Industrial Organization}, 79:102712,
  2021.

\bibitem[Deng et~al.(2025)]{deng2025exploring}
S.~Deng, M.~Schiffer, and M.~Bichler.
\newblock Exploring competitive and collusive behaviors in algorithmic
  pricing with deep reinforcement learning.
\newblock \emph{arXiv preprint arXiv:2503.11270}, 2025.

\bibitem[Fudenberg et~al.(1994)]{fudenberg1994folk}
D.~Fudenberg, D.~K.~Levine, and E.~Maskin.
\newblock The folk theorem with imperfect public information.
\newblock \emph{Econometrica}, 62(5):997--1039, 1994.

\bibitem[Green and Porter(1984)]{green1984noncooperative}
E.~J.~Green and R.~H.~Porter.
\newblock Noncooperative collusion under imperfect price information.
\newblock \emph{Econometrica}, 52(1):87--100, 1984.

\bibitem[Schulman et~al.(2017)]{schulman2017proximal}
J.~Schulman, F.~Wolski, P.~Dhariwal, A.~Radford, and O.~Klimov.
\newblock Proximal policy optimization algorithms.
\newblock \emph{arXiv preprint arXiv:1707.06347}, 2017.

\bibitem[Vehtari et~al.(2021)]{vehtari2021rank}
A.~Vehtari, A.~Gelman, D.~Simpson, B.~Carpenter, and P.-C.~B\"urkner.
\newblock Rank-normalization, folding, and localization: An improved
  $\hat{R}$ for assessing convergence of {MCMC}.
\newblock \emph{Bayesian Analysis}, 16(2):667--718, 2021.

\end{thebibliography}

%% ============================================================
%% AI Use Disclosure
%% ============================================================
\section*{AI Use Disclosure}
Generative AI was used to assist with manuscript drafting and revision, mathematical derivation checks, implementation review, postprocessing of archived experimental outputs, and figure preparation. All numerical results were computed from the archived simulator runs rather than generated hypothetically, and the authors retain full responsibility for all assumptions, derivations, code, references, and claims.

\end{document}
